\documentclass[11pt,a4paper]{article}

% Working notes for the Neural Particle Method project.
% Scrap document: ideas as discussed, not paper prose. 2026-08-31.

\usepackage{amsmath, amssymb, amsthm, mathtools}
\usepackage{bm}
\usepackage[numbers,sort&compress]{natbib}
\usepackage[colorlinks=true,linkcolor=blue,citecolor=blue,urlcolor=blue]{hyperref}
\usepackage[capitalise,noabbrev]{cleveref}
\usepackage{microtype}
\usepackage{parskip}
\setlength{\parindent}{0pt}
\setlength{\parskip}{6pt}
\usepackage[margin=2.5cm]{geometry}
\usepackage{xcolor}

\theoremstyle{plain}
\newtheorem{theorem}{Theorem}[section]
\newtheorem{lemma}[theorem]{Lemma}
\newtheorem{proposition}[theorem]{Proposition}
\newtheorem{corollary}[theorem]{Corollary}

\theoremstyle{definition}
\newtheorem{definition}[theorem]{Definition}
\newtheorem{assumption}{Assumption}

\theoremstyle{remark}
\newtheorem{remark}[theorem]{Remark}

\newcommand{\inlinetodo}[1]{\textcolor{red}{\textbf{[TODO: #1]}}}

\newcommand{\E}{\mathbb{E}}
\newcommand{\Pm}{\mathbb{P}}
\newcommand{\Qm}{\mathbb{Q}}
\newcommand{\R}{\mathbb{R}}
\newcommand{\norm}[1]{\left\lVert #1 \right\rVert}
\newcommand{\abs}[1]{\left\lvert #1 \right\rvert}
\newcommand{\sigDup}{\sigma_{\mathrm{Dup}}}
\newcommand{\Lev}{L}
\newcommand{\Phimap}{\Phi}
\newcommand{\Phal}{\Phi_\alpha}
\newcommand{\Id}{\mathrm{Id}}

\title{Neural Particle Method: Convergence Notes for the Implicit Scheme}
\author{Osian Shelley}
\date{31 August 2026}

\begin{document}

\maketitle

\begin{abstract}
  Scrap document holding the convergence analysis sketch for the implicit
  scheme, split out from the main working notes
  (\texttt{notes.tex}).
\end{abstract}

\section{Context}\label{sec:context}

Recalled from the main notes: the LSV model
$dX_t/X_t = \Lev(t,X_t)\sqrt{V_t}\,dW_t$ is calibrated when the leverage
satisfies $\Lev = \sigDup/\sqrt{\E[V_t \mid X_t = \cdot\,]}$, the conditional
expectation being taken under the law generated by $\Lev$ itself. A
calibrated model is a fixed point of the self-consistency map
\begin{equation}\label{eq:phimap}
  \Phimap\colon \Lev \longmapsto
  \frac{\sigDup}{\sqrt{\E^{\Lev}[V_t \mid X_t = \cdot\,]}},
\end{equation}
and the implicit scheme iterates the damped update
\begin{equation}\label{eq:damped}
  \Lev^{n+1} = (1 - \alpha)\,\Lev^{n} + \alpha\,\Phimap(\Lev^{n}),
  \qquad \alpha \in (0, 1].
\end{equation}
The explicit per-slice scheme contains no fixed point; everything below
concerns the implicit scheme only.

\section{Convergence analysis sketch}\label{sec:theory}

\subsection{Damping algebra}

Write $\Phal \coloneqq (1-\alpha)\Id + \alpha\Phimap$ and linearise at a fixed
point $\Lev^*$. Each spectral value $\lambda$ of $D\Phimap(\Lev^*)$ is mapped to
\begin{equation}\label{eq:mumap}
  \mu = (1-\alpha) + \alpha\lambda = 1 + \alpha(\lambda - 1).
\end{equation}
One has $\abs{\mu} < 1$ for some $\alpha > 0$ if and only if
$\operatorname{Re}\lambda < 1$, with the per-eigenvalue threshold
\begin{equation}\label{eq:threshold}
  \alpha < \frac{2\,(1 - \operatorname{Re}\lambda)}{\abs{1 - \lambda}^2},
\end{equation}
and the overall threshold is the infimum of \cref{eq:threshold} over the
spectrum. Two consequences worth recording.
\begin{enumerate}
  \item Damping cannot rescue any spectral value with
    $\operatorname{Re}\lambda \geq 1$. For real $\lambda > 1$,
    $\mu = 1 + \alpha(\lambda - 1) > 1$ for every $\alpha > 0$. The structural
    hypothesis below is therefore the real content of the convergence claim.
  \item Damping never changes the fixed-point set,
    $\Phal(\Lev) = \Lev \iff \Phimap(\Lev) = \Lev$. Damping buys convergence,
    never existence. Existence must come from a separate argument.
\end{enumerate}

\begin{assumption}[Spectral condition]\label{ass:spectral}
  The spectrum of $D\Phimap(\Lev^*)$ is contained in
  $\{z \in \mathbb{C} : \operatorname{Re} z < 1\}$.
\end{assumption}

Heuristic support for \cref{ass:spectral}: the oscillatory overshoot observed
in undamped iterations of the classical calibration loop suggests spectral
values with negative real part, which is exactly the regime where damping helps
most. \inlinetodo{Explore estimating the dominant eigenvalue numerically, e.g.\
power iteration on the calibration loop; would turn the assumption into a
measurement.}

\subsection{Existence versus local convergence: two separate routes}

The following two arguments must not be conflated.

\emph{Ostrowski (spectral radius) route.} Assuming a fixed point $\Lev^*$
exists and $\Phimap$ is Fr\'echet differentiable at $\Lev^*$, the condition
$\rho\bigl(D\Phal(\Lev^*)\bigr) < 1$ yields local linear convergence of
\cref{eq:damped} from a sufficiently good initial guess, via the Gelfand
equivalent-norm construction. In function space one must speak of the spectrum
rather than eigenvalues. This is a local convergence theorem and gives no
existence.

\emph{Banach route.} Existence and uniqueness require $\Phal$ to be a
contraction on a complete metric space that it maps into itself. At the linear
level in a Hilbert norm the sharp condition concerns the \emph{numerical range}
of $D\Phimap$, which contains the spectrum and is in general strictly larger:
one needs the numerical range in $\{\operatorname{Re} z < 1\}$ together with a
bound on $\norm{D\Phimap}$. Stronger hypothesis, global conclusion (existence,
uniqueness, geometric convergence).

\emph{Honest fallback structure.} Establish existence by a short-horizon
contraction argument (Lipschitz constant of $\Phimap$ of order $CT$, matching
the short-time results in the literature), or by Schauder compactness for the
discretised problem (existence without uniqueness); then apply the Ostrowski
route for local convergence of the damped iteration.

\subsection{Fr\'echet differentiability of $\Phimap$: the hard step}

The map $\Phimap$ factors as
$\Lev \mapsto \text{law of } (X, V) \mapsto \E[V_t \mid X_t = \cdot\,]
\mapsto \sigDup / \sqrt{\cdot}$, and with joint density $p$,
\[
  \E[V_t \mid X_t = x] = \frac{\int v\, p(x, v)\, dv}{\int p(x, v)\, dv},
\]
so differentiating $\Phimap$ in $\Lev$ means differentiating a ratio of
densities with respect to an SDE coefficient. Required ingredients:
\begin{enumerate}
  \item existence and smoothness of the joint density of $(X_t, V_t)$;
  \item a strictly positive lower bound on the $X$-marginal on the region
    where $\Lev$ is evaluated;
  \item differentiability of $\Lev \mapsto p^{\Lev}$, by Malliavin calculus or
    parabolic PDE perturbation theory.
\end{enumerate}
The denominator bound in item 2 is the tail-degeneracy issue from practice
reappearing in the theory. Capping or truncating the spatial domain is not an
engineering hack; it is what makes $\Phimap$ differentiable on a manageable
set. This alignment deserves a remark in the paper.

\subsection{The implementable theorem}

The algorithm iterates a perturbed map $\widehat{\Phimap} = \Phimap +
\varepsilon$, where $\varepsilon$ collects the neural regression error and the
finite-particle error. Target statement: under \cref{ass:spectral} and for
$\alpha$ below the threshold \cref{eq:threshold}, the damped iteration on
$\widehat{\Phimap}$ converges to an $\varepsilon/(1-r)$-neighbourhood of
$\Lev^*$, where $r < 1$ is the contraction factor of the linearised damped map.
This has the same structure as the deep BSDE error analyses, and it is the
statement the experiments can verify.

\subsection{State of the well-posedness literature}

Well-posedness of the calibrated McKean--Vlasov equation is open in full
generality. Milestones: short-time existence for the associated nonlinear PIDE
(Abergel--Tachet 2010); existence for finite-state volatility (Jourdain--Zhou);
inversion of the Markovian projection and stationary solutions under
restrictions (Lacker--Shkolnikov--Zhang 2019); existence under non-regular
coefficients with propagation of chaos (Djete 2022); strong
existence-uniqueness for finite-state volatility with propagation of chaos
(Mustapha 2024). Continuous-state Heston-type LSV remains unresolved without
regularisation. Any fixed-point claim should therefore be made for the
discretised or regularised map, or under \cref{ass:spectral} as a named
hypothesis, and not in a generality exceeding these results. See
\texttt{papers/READING\_LIST.md} for the annotated list.

\end{document}
